% ============================================================
%  MODULE VIII — Market Microstructure
% ============================================================

\chapter{Market Microstructure}

\begin{notebox}
Market microstructure is the study of how markets actually function at a mechanical level —
order books, liquidity, price formation, and the behaviour of market participants at
fine-grained timescales. Understanding this gives traders an edge in execution and
helps explain ``why prices do what they do.''
\end{notebox}

% ────────────────────────────────────────────────────────────
\section{The Order Book}

The \textbf{order book} (Level 2 data) is a real-time record of all outstanding buy
and sell limit orders at each price level.

\begin{center}
\begin{tabular}{p{3.5cm} >{\bfseries}p{3cm} p{3.5cm}}
\toprule
Bid (Buy) Size & Price & Ask (Sell) Size \\
\midrule
500 & 99.97 & \\
1200 & 99.98 & \\
3000 & 99.99 & \\
\multicolumn{3}{c}{\textbf{--- SPREAD ---}} \\
& 100.01 & 2500 \\
& 100.02 & 1100 \\
& 100.03 & 800 \\
\bottomrule
\end{tabular}
\end{center}

\begin{itemize}
  \item \textbf{Bid} — highest price a buyer is willing to pay (\$99.99)
  \item \textbf{Ask} (Offer) — lowest price a seller is willing to accept (\$100.01)
  \item \textbf{Spread} — difference = \$0.02 in this case
  \item \textbf{Mid price} — (Bid + Ask) / 2 = \$100.00
  \item \textbf{Depth} — the size available at each price level
\end{itemize}

% ────────────────────────────────────────────────────────────
\section{Liquidity and Market Impact}

\begin{keyterm}{Liquidity}
The ease with which a security can be bought or sold without meaningfully affecting its price.
Measured by: spread (tighter = more liquid), depth (more size = more liquid), and
market impact (large orders moving the price).
\end{keyterm}

A \textbf{market order} eats through the order book. A large market buy order
``walks up the book'', consuming all available ask orders until filled — each successive fill
occurs at a higher price. This is \textbf{market impact / slippage}.

\begin{notebox}
For small retail traders, slippage is negligible on liquid large-cap stocks.
For institutional traders executing millions of shares, market impact is a major cost.
This is why institutions use algorithmic execution (VWAP, TWAP, Implementation Shortfall).
\end{notebox}

% ────────────────────────────────────────────────────────────
\section{Bid-Ask Spread as a Cost}

Every time you cross the spread, you pay a cost:

\begin{formulabox}
\[
  \text{Spread Cost (round trip)} = \frac{\text{Ask} - \text{Bid}}{\text{Mid}} \times 2 \times \text{Position Value}
\]
\end{formulabox}

For a stock with a \$0.05 spread at \$50 mid, round-trip spread cost = 0.2\%.
At \$100{,}000 position size: \$200 per round trip — before any other cost.

% ────────────────────────────────────────────────────────────
\section{Types of Orders and the Book}

\begin{itemize}
  \item \textbf{Limit order} — added to the book; provides liquidity. Often earns a
        \textbf{maker rebate} on exchanges.
  \item \textbf{Market order} — removes from the book; takes liquidity. Pays a
        \textbf{taker fee}.
  \item \textbf{Iceberg order} — large order with only part visible; hides intent.
  \item \textbf{Stop order} — hidden from the book; triggered and becomes market order at stop price.
\end{itemize}

% ────────────────────────────────────────────────────────────
\section{Price Discovery}

\begin{keyterm}{Price Discovery}
The process by which markets aggregate information from all participants into a single
observed price. More informative prices = more efficient allocation of capital.
\end{keyterm}

Price discovery happens continuously during market hours. \textbf{Opening auctions} and
\textbf{closing auctions} are special mechanisms that batch orders to set the open and
close price — the most liquid moments of the trading day.

\subsection{Pre-Market and After-Hours Trading}
Trading outside regular hours (US: 9:30am--4:00pm ET). Characteristics:
\begin{itemize}
  \item Much lower liquidity and volume
  \item Wider spreads
  \item More volatile price action
  \item Earnings releases often happen after hours or pre-market
\end{itemize}

% ────────────────────────────────────────────────────────────
\section{Volume Analysis}

Volume is the number of shares/contracts traded in a period.
It is the \textbf{second most important data point after price}.

\begin{itemize}
  \item \textbf{High volume breakout} — conviction; more participants validating the move.
  \item \textbf{Low volume breakout} — suspect; may fail or reverse.
  \item \textbf{Climax volume} — exhaustion spike; potential reversal signal.
  \item \textbf{Volume dry-up} — lack of sellers (in uptrend) or buyers (in downtrend);
        often precedes continuation.
\end{itemize}

\subsection{VWAP (Volume-Weighted Average Price)}

\begin{formulabox}
\[
  VWAP = \frac{\sum (Price_i \times Volume_i)}{\sum Volume_i}
\]
\end{formulabox}

VWAP is the average price weighted by volume. Institutional traders use it as a benchmark.
Trading above VWAP = bullish intraday; below = bearish intraday.

% ────────────────────────────────────────────────────────────
\section{Market Regimes and Volatility}

\subsection{VIX — The Fear Index}
The \textbf{CBOE Volatility Index (VIX)} measures the market's expectation of
30-day S\&P 500 volatility, derived from options pricing.

\begin{itemize}
  \item VIX $<$ 15: Low volatility, complacent market
  \item VIX 15--25: Normal market
  \item VIX 25--40: Elevated fear
  \item VIX $>$ 40: Crisis, extreme fear (March 2020 hit 85)
\end{itemize}

\begin{notebox}
High VIX environments have \textbf{wider spreads}, \textbf{faster moves}, and
\textbf{larger gaps}. Position sizes should be \emph{reduced} in high volatility
regimes to maintain constant dollar risk.
\end{notebox}

% ────────────────────────────────────────────────────────────
\section{Order Flow and Tape Reading}

\textbf{Order flow} refers to the \emph{real-time stream of transactions} executed in the market.
Advanced traders read the tape (time and sales) to detect:

\begin{itemize}
  \item Aggressive buying (large market orders lifting the ask)
  \item Aggressive selling (large orders hitting the bid)
  \item Absorption (large hidden orders absorbing flow without price movement)
  \item \textbf{Stop runs} — price briefly moves to known stop loss zones, triggers stops,
        then reverses. Common manipulation tactic by larger players.
\end{itemize}

\begin{warningbox}
Order flow analysis has a steep learning curve. Do not attempt it until you have mastered
basic chart reading, risk management, and have consistent results on higher timeframes first.
\end{warningbox}
